\section{Caracterización del proyecto}

\subsection{Objetivo general}
Diseñar, implementar y validar un lazo de control para un convertidor DC-DC bidireccional en configuración interleaved, optimizando la reducción de la inductancia y el rizado de corriente, utilizando una plataforma HIL para su validación en tiempo real.
\subsection{Objetivos específicos}
\begin{itemize}
    \item Justificar y demostrar que la configuración interleaved permite reducir a la mitad la frecuencia efectiva de operación del convertidor.

    \item Analizar cómo la reducción en la frecuencia de conmutación impacta en el valor requerido de la inductancia, permitiendo utilizar inductores de menor tamaño y costo sin comprometer el rendimiento.

    \item Evaluar la disminución del rizado de corriente mediante la combinación de corrientes desfasadas de las fases interleaved, mejorando la calidad de la señal y reduciendo las pérdidas térmicas.

    \item Desarrollar e implementar el lazo de control que sincronice y coordine ambas fases para maximizar los beneficios de la configuración interleaved en términos de reducción de rizado y eficiencia.
\end{itemize}







\subsection{Alcance}
